%File: simvla_tds_aaai2027.tex
% =============================================================================
% AAAI-2027 anonymous submission.
% Core contribution: Transition-Aware Dynamic Sampling (TDS) for sample-efficient
% Vision-Language-Action learning, built on the SimVLA base model.
%
% TODO BEFORE SUBMISSION (search "TODO"):
%   * Replace the fbox figure placeholders with the generated PDFs
%     (architecture, efficiency curve = main_suite_curves.png, Exp-0 scatter =
%     density_vs_sr.png, per-task appendix grids).
%   * Fill the ablation table rows marked "--" once the control runs finish
%     (measured-difficulty / uniform-interleaved / inverted / no-clip).
%   * Confirm the published-baseline numbers against the cited sources.
% Compile: pdflatex -> bibtex -> pdflatex x2, with aaai2027.sty/.bst + references.bib.
% =============================================================================
\documentclass[letterpaper]{article} % DO NOT CHANGE THIS
\usepackage[submission]{aaai2027}  % DO NOT CHANGE THIS
\usepackage[hyphens]{url}  % DO NOT CHANGE THIS
\usepackage{graphicx} % DO NOT CHANGE THIS
\urlstyle{rm} % DO NOT CHANGE THIS
\def\UrlFont{\rm}  % DO NOT CHANGE THIS
\usepackage{natbib}  % DO NOT CHANGE THIS AND DO NOT ADD ANY OPTIONS TO IT
\usepackage{caption} % DO NOT CHANGE THIS AND DO NOT ADD ANY OPTIONS TO IT
\frenchspacing  % DO NOT CHANGE THIS

\usepackage{algorithm}
\usepackage{algorithmic}
\usepackage{booktabs}
\usepackage{amsmath}
\usepackage{amssymb}
\usepackage{multirow}

\pdfinfo{
/TemplateVersion (2027.1)
}

\setcounter{secnumdepth}{0}

% ---- convenience macros ----
\newcommand{\tds}{TDS}
\newcommand{\dk}{d_{k}}
\newcommand{\pk}{p_{k}}
\newcommand{\na}{\textbf{[TODO]}}

\title{Transition-Aware Dynamic Sampling for Sample-Efficient\\Vision-Language-Action Learning}
\author{
    Anonymous Submission
}
\affiliations{
    Anonymous Submission
}

\begin{document}

\maketitle

\begin{abstract}
Vision-language-action (VLA) models are almost always trained by minimizing a uniform behavioral-cloning objective over uniformly sampled demonstrations, implicitly assuming that every task in the training mixture is equally informative. Manipulation difficulty, however, is highly non-uniform, and we observe that it is already encoded in the demonstrations themselves: the density of gripper open/close events, the fraction of low-speed fine-alignment segments, and the trajectory length jointly characterize how hard a task is to imitate. We introduce \emph{Transition-Aware Dynamic Sampling} (\tds{}), a training strategy that scores per-task difficulty \emph{offline} from this transition structure and dynamically reweights how often each task is drawn during training, concentrating optimization on the hard, under-learned tasks. \tds{} requires no reference model, no periodic policy evaluation, and adds no measurable training overhead, in contrast to loss- or success-based data-mixing methods that must repeatedly query the model. Built on the compact SimVLA base model and evaluated on the LIBERO benchmark, \tds{} improves average success rate by $+27.9$ points at a fixed $40$k-step budget (same-machine controlled comparison) and reaches the uniform baseline's $100$k-step performance in roughly half the steps, while remaining ahead at convergence. An offline-computed difficulty score also correlates significantly with measured per-task failure ($\rho=-0.32$, $p<0.05$), supporting the premise that transition structure is a model-free proxy for task difficulty.
\end{abstract}

\section{Introduction}

Vision-language-action (VLA) models map raw visual observations and natural-language instructions directly to robot actions, and have advanced quickly by coupling pretrained vision-language backbones with expressive action decoders \citep{brohan2023rt2,kim2024openvla,octo2024,black2024pi0}. The prevailing training recipe is behavioral cloning under a \emph{uniform} objective: each timestep of each demonstration of each task contributes equally to the loss, and minibatches are sampled uniformly from the data mixture.

This uniformity sits uneasily with the structure of manipulation. Across a task suite, difficulty varies by a large margin---some tasks are mastered almost immediately while long-horizon, multi-stage tasks lag by tens of success-rate points \citep{liu2023libero,kim2024openvla}. Under uniform sampling, a model spends the same optimization budget rehearsing already-solved easy tasks as it does on the hard tasks that actually determine average performance. The result is poor \emph{sample efficiency}: the policy reaches a given success rate only after far more gradient steps than necessary.

A natural remedy is to sample hard tasks more often. The dominant way to decide \emph{which} tasks are hard is to \emph{ask the model}: train a reference or proxy model and read off difficulty from its loss, temporal-difference error, or rollout success \citep{schaul2016prioritized,lin2017focal,hejna2024remix,xie2023doremi}. This couples the difficulty estimate to the very model being trained, changes when the backbone changes, and introduces an extra evaluation loop with its own cost and hyperparameters.

We take the opposite view: \emph{difficulty is a property of the demonstrations, observable before training begins}. The moment a demonstration is recorded, its action signal already reveals how many sub-task events it contains (gripper open/close events), how much of it demands slow, low-tolerance alignment (low-speed plateaus), and how much room it leaves for error accumulation (trajectory length). All three are deterministic functions of the action sequence, computable in a single offline pass, independent of any model.

Building on this observation we propose \emph{Transition-Aware Dynamic Sampling} (\tds{}). \tds{} (i) computes a per-task transition-density difficulty score offline from the demonstrations, (ii) converts it into a sampling distribution over tasks via a temperature-controlled transform with a uniform floor, and (iii) drives a streaming, two-level sampler that dynamically realizes this distribution during training---all while leaving the base model, the loss, and the optimizer untouched. Because the difficulty signal is data-intrinsic, \tds{} adds no training-time evaluation and is identical for any backbone.

Our contributions are:
\begin{itemize}
\item \textbf{The TDS mechanism}: a model-free, zero-training-cost characterization of per-task manipulation difficulty from action transition structure, and a temperature-with-floor reweighting that turns it into a task-sampling distribution.
\item \textbf{Integration into SimVLA}: a streaming two-level sampler that realizes the target distribution exactly over an infinite iterable dataset, modifying only the data path of an otherwise standard flow-matching VLA.
\item \textbf{Improved sample efficiency}: on LIBERO, \tds{} yields $+27.9$ average success rate at a fixed low budget and $\sim\!2\times$ fewer steps to match the uniform baseline (same-machine controlled comparison), with the offline difficulty score significantly predicting measured per-task failure.
\end{itemize}

\begin{figure}[t]
\centering
\fbox{\parbox{0.92\columnwidth}{\centering \vspace{2.2cm} TODO: SimVLA $+$ TDS overview. Offline difficulty $\dk$ from transition structure (gripper events, low-speed plateaus, length) $\rightarrow$ task sampling weights $\pk$ $\rightarrow$ streaming two-level sampler feeding the SmolVLM$+$DiT flow-matching policy. \vspace{2.2cm}}}
\caption{Overview of SimVLA with Transition-Aware Dynamic Sampling. Difficulty is read from the demonstrations before training; only the sampler changes.}
\label{fig:overview}
\end{figure}

\section{Related Work}

\paragraph{Vision-language-action models.}
Large VLA models attach action decoders to pretrained vision-language backbones, either by discretizing actions into tokens \citep{brohan2022rt1,brohan2023rt2,kim2024openvla} or by attaching continuous diffusion/flow heads \citep{octo2024,black2024pi0,chi2023diffusion}. Recent work also optimizes the fine-tuning recipe itself \citep{kim2025oft}. These efforts keep the training objective and sampling uniform over tasks; \tds{} is complementary and modifies only how tasks are sampled.

\paragraph{Imitation learning and behavioral cloning.}
Action chunking with diffusion or flow matching has become a standard, effective recipe for visuomotor imitation \citep{chi2023diffusion,zhao2023aloha,lipman2023flow}. Our base model follows this paradigm; \tds{} is agnostic to the specific decoder.

\paragraph{Curriculum learning and data sampling.}
Emphasizing hard examples has a long history: prioritized replay \citep{schaul2016prioritized}, focal loss \citep{lin2017focal}, online hard-example mining \citep{shrivastava2016ohem}, and curriculum learning \citep{bengio2009curriculum}. For data mixtures, Re-Mix \citep{hejna2024remix} learns domain weights via distributionally robust optimization over a reference model, and DoReMi \citep{xie2023doremi} trains a proxy model to set mixture proportions, reaching baseline accuracy in $2.6\times$ fewer steps. All of these derive difficulty from a \emph{model}. \tds{} differs in kind: its weights are a property of the demonstrations, computed offline, identical for any backbone, and free of an evaluation loop.

\paragraph{Transition modeling.}
Most VLA training treats all transitions within a trajectory as equally important. We instead use the \emph{aggregate} transition structure of a task's demonstrations as a difficulty signal for sampling, rather than modeling individual transitions; this aggregate is exactly what is missing from uniform task sampling.

\section{Method}

\subsection{Problem Formulation}
A demonstration is a trajectory $\tau=\{(o_t,a_t)\}_{t=1}^{T}$ of observations $o_t$ (multi-view images, language instruction, proprioception) and actions $a_t\in\mathbb{R}^{d}$, inducing state transitions $s_t\!\to\!s_{t+1}$. The training set is partitioned into $K$ tasks $\{\mathcal{D}_k\}_{k=1}^{K}$. A policy $\pi_\theta$ is trained by behavioral cloning over action chunks $a_{t:t+H}$. Writing $p$ for the distribution from which a task is drawn at each training step, the objective is
\begin{equation}
\min_\theta\ \mathbb{E}_{k\sim p}\,\mathbb{E}_{(o,a)\sim\mathcal{D}_k}\big[\,\mathcal{L}\!\left(\pi_\theta(o),a\right)\big].
\label{eq:objective}
\end{equation}
The standard recipe sets $p$ uniform, $p_k=1/K$. \tds{} replaces $p$ with a difficulty-aware distribution $\pk$ derived offline from the demonstrations, leaving $\mathcal{L}$ and $\theta$'s optimization unchanged.

\subsection{SimVLA Base Model}
SimVLA is a compact flow-matching VLA. A pretrained SmolVLM vision-language backbone \citep{marafioti2025smolvlm} encodes the multi-view images and the instruction into conditioning features; a DiT-style action transformer \citep{peebles2023dit} consumes these features together with proprioception and a noised action chunk, and is trained with conditional flow matching \citep{lipman2023flow} to predict the chunk's velocity field. At inference the head integrates the velocity to produce an $H$-step action chunk. SimVLA is intentionally simple and serves only as the substrate on which \tds{} operates; \tds{} makes no assumption about the decoder beyond the loss in Eq.~\eqref{eq:objective}.

\subsection{TDS: Transition-Aware Dynamic Sampling}

\paragraph{Motivation.}
Not all tasks are equally informative to train on. Tasks that are already solved contribute little gradient signal, whereas hard tasks---those with many contact/grasp events, long low-tolerance alignment phases, and long horizons---remain under-learned under uniform sampling. Crucially, these properties are visible in the demonstrations \emph{before} training, so we can prioritize hard tasks without ever querying the model.

\paragraph{The transition signal.}
For an episode with actions $a_{1:T}$ we define a per-dimension-normalized change rate
\begin{equation}
c_t=\frac{1}{|\mathcal{J}|}\sum_{j\in\mathcal{J}}\frac{\lvert a_{t,j}-a_{t-1,j}\rvert}{\max_{t'}\lvert a_{t',j}-a_{t'-1,j}\rvert+\varepsilon},
\label{eq:changerate}
\end{equation}
where $\mathcal{J}$ excludes the gripper dimension and $c$ is renormalized to $[0,1]$ within the episode. From this we extract three task statistics, each averaged over the task's demonstrations:
(1) \textbf{gripper events} $E$: the number of sign flips of the (median-centered) gripper command, robust to $0/1$ and $-1/1$ encodings, counting grasp/release sub-tasks;
(2) \textbf{plateau ratio} $P$: the fraction of steps inside low-speed plateaus, i.e.\ maximal runs of length $\geq L_p$ with $c_t<\theta_p$, capturing fine-alignment phases with narrow tolerance;
(3) \textbf{length} $L$: the mean trajectory length, a proxy for error-accumulation horizon. We use $\theta_p{=}0.25$, $L_p{=}3$.

\paragraph{Difficulty score.}
Because the three statistics have incommensurate units, each is standardized across tasks and combined into a difficulty score
\begin{equation}
\dk=\alpha\, z(E_k)+\beta\, z(P_k)+\gamma\, z(L_k),
\label{eq:dk}
\end{equation}
with $z(\cdot)$ the across-task $z$-score. On LIBERO we use the two-component form $\alpha{=}0,\ \beta{=}\gamma{=}\tfrac12$: an analysis (below) shows the gripper-event term carries no discriminative signal on this benchmark, so it is dropped. Note $\dk$ is centered at $0$; negative values denote easier-than-average tasks.

\paragraph{Sampling weights.}
We map $\dk$ to a sampling distribution with a temperature transform and a uniform floor:
\begin{equation}
\tilde{p}_k\propto\big(\dk-\textstyle\min_{k'}d_{k'}+1\big)^{1/T},\qquad
\pk\propto\max\!\big(\tilde{p}_k,\ \rho/K\big),
\label{eq:pk}
\end{equation}
with temperature $T{=}2$ and floor $\rho{=}0.5$. The shift makes the base positive so the fractional power is well defined; $T$ controls how aggressively hard tasks are up-weighted ($T\!\to\!\infty$ recovers uniform sampling). The floor guarantees every task at least half its uniform share, preventing catastrophic forgetting of easy tasks.

\paragraph{Dynamic streaming sampler.}
SimVLA streams from an infinite iterable dataset, so a map-style weighted sampler does not apply. Instead we use a two-level scheme: at each draw we first sample a task $k\sim\pk$, then take the next sample from that task's persistent episode stream. Because the weighted draw occurs once per \emph{sample}, the empirical fraction of training samples from task $k$ converges exactly to $\pk$, independent of per-task demo counts or episode lengths; a monitor verifies the empirical distribution against $\pk$ early in training. Setting $\pk$ uniform recovers the original training stream up to sampling order. The training objective is unchanged from Eq.~\eqref{eq:objective} with $p=\pk$.

\begin{algorithm}[tb]
\caption{Transition-Aware Dynamic Sampling (\tds{})}
\label{alg:tds}
\textbf{Input}: demos $\{\mathcal{D}_k\}_{k=1}^{K}$, base policy $\pi_\theta$, temperature $T$, floor $\rho$\\
\textbf{Offline (once, CPU)}: for each task $k$ compute $E_k,P_k,L_k$; $\dk$ (Eq.~\ref{eq:dk}); $\pk$ (Eq.~\ref{eq:pk})\\
\textbf{Output}: trained $\pi_\theta$
\begin{algorithmic}[1]
\WHILE{training}
\STATE sample task $k\sim\pk$
\STATE draw next chunk $(o,a)$ from task $k$'s stream
\STATE update $\theta$ on $\mathcal{L}(\pi_\theta(o),a)$
\ENDWHILE
\end{algorithmic}
\end{algorithm}

\subsection{Extensions (Future Work)}
The transition signal of Eq.~\eqref{eq:changerate} is defined at the timestep level and could in principle reweight the per-step loss or gate execution granularity. We focus this paper on the task-sampling instantiation, which is the change responsible for the gains reported below, and leave step- and execution-level uses to future work.

\section{Experiments}

\subsection{Setup}
\paragraph{Benchmark.}
We evaluate on LIBERO \citep{liu2023libero}, four suites of ten tasks each: \textsc{Spatial} (layout generalization), \textsc{Object} (object generalization), \textsc{Goal} (goal generalization), and \textsc{Long} (long-horizon, multi-stage). We report per-suite and average success rate (SR, \%) over $20$ trials per task ($200$ rollouts per suite).
\paragraph{Implementation.}
SimVLA uses SmolVLM-500M-Instruct \citep{marafioti2025smolvlm} with a DiT action transformer ($H{=}10$, $384^2$ images, two views). \tds{} and the uniform baseline share an identical recipe (batch size, learning rate, schedule, seed); the \emph{only} difference is the task-sampling distribution. \tds{} statistics are computed once from the raw demonstrations on CPU in minutes. To remove cross-machine variance, the headline $40$k/$100$k comparisons are evaluated for both models on the same machine.

\subsection{Main Results: Sample Efficiency}
Table~\ref{tab:main} reports the controlled, same-machine comparison at a low budget ($40$k steps) and at convergence ($100$k steps). At $40$k, \tds{} improves average SR from $61.9$ to $89.75$ ($+27.9$), with the largest gains on the hard \textsc{Goal} ($+53.5$) and \textsc{Long} ($+30.5$) suites---exactly the tasks \tds{} up-weights. At $100$k the suites approach saturation, yet \tds{} still leads by $+4.0$, indicating it does not merely accelerate training but also reaches a higher plateau, with gains concentrated on the hard suites. This $+27.9$ is measured against a naive uniform baseline; the ablation in Table~\ref{tab:ablation} decomposes it into a sampler-mechanism component (batch decorrelation) and the transition-density weighting itself, and we attribute the incremental gain over a well-mixed uniform sampler to the weighting.

\begin{table*}[t]
\centering
\begin{tabular}{llccccc}
\toprule
Budget & Method & \textsc{Spatial} & \textsc{Object} & \textsc{Goal} & \textsc{Long} & Avg \\
\midrule
\multirow{2}{*}{$40$k steps}
& SimVLA (uniform) & 73.0 & 95.5 & 38.5 & 40.5 & 61.9 \\
& SimVLA $+$ \tds{} (ours) & \textbf{97.5} & \textbf{98.5} & \textbf{92.0} & \textbf{71.0} & \textbf{89.75} \\
\cmidrule(lr){2-7}
& $\Delta$ & $+24.5$ & $+3.0$ & $+53.5$ & $+30.5$ & $+27.9$ \\
\midrule
\multirow{2}{*}{$100$k steps}
& SimVLA (uniform) & 92.5 & 96.0 & 90.5 & 87.5 & 91.6 \\
& SimVLA $+$ \tds{} (ours) & \textbf{97.5} & \textbf{98.5} & \textbf{96.0} & \textbf{90.5} & \textbf{95.6} \\
\cmidrule(lr){2-7}
& $\Delta$ & $+5.0$ & $+2.5$ & $+5.5$ & $+3.0$ & $+4.0$ \\
\bottomrule
\end{tabular}
\caption{\textbf{Main result: same-machine controlled comparison on LIBERO} (SR \%, $20$ trials/task). The only difference between the two rows in each block is the task-sampling distribution. \tds{} yields a large gain at low budget that narrows but stays positive at convergence.}
\label{tab:main}
\end{table*}

Table~\ref{tab:curve} gives the full efficiency curve. \tds{} at $20$k steps ($75.6$) already exceeds the uniform baseline at $40$k ($61.9$); \tds{} at $60$k ($95.0$) exceeds uniform at $100$k ($91.6$). Thus \tds{} reaches the uniform baseline's converged success rate in roughly half the training steps. (Curve points for \tds{} at $40$/$60$/$80$k are evaluated on a second machine and are mildly conservative relative to the same-machine anchors in Table~\ref{tab:main}; the qualitative dominance is unaffected.)

\begin{table}[t]
\centering
\setlength{\tabcolsep}{1.2mm}
\begin{tabular}{lccccc}
\toprule
Steps (k) & 20 & 40 & 60 & 80 & 100 \\
\midrule
Uniform & 35.1 & 61.9 & 85.4 & 81.5 & 91.6 \\
\tds{} (ours) & \textbf{75.6} & \textbf{86.9} & \textbf{95.0} & \textbf{93.9} & \textbf{95.6} \\
\midrule
Gap & $+40.5$ & $+25.0$ & $+9.6$ & $+12.4$ & $+4.0$ \\
\bottomrule
\end{tabular}
\caption{\textbf{Efficiency curve} (average SR \% vs.\ training steps). The gap is largest early and stays positive throughout. See Fig.~\ref{fig:curve}.}
\label{tab:curve}
\end{table}

\begin{figure}[t]
\centering
\fbox{\parbox{0.9\columnwidth}{\centering \vspace{2.4cm} TODO: efficiency-curve figure (avg SR vs.\ steps, uniform vs.\ \tds{}, per-suite panels) --- insert \texttt{main\_suite\_curves.png}. \vspace{2.4cm}}}
\caption{Average and per-suite success rate vs.\ training steps. \tds{} (ours) dominates the uniform baseline across budgets and suites, reaching the baseline's $100$k performance in $\sim\!2\times$ fewer steps.}
\label{fig:curve}
\end{figure}

\paragraph{Context with published methods.}
For reference, on LIBERO the reported averages are $72.4$ (Diffusion Policy), $75.1$ (Octo), $76.5$ (OpenVLA), and $97.1$ (OpenVLA-OFT) \citep{chi2023diffusion,octo2024,kim2024openvla,kim2025oft}, obtained under each method's own backbone and training budget. These contextualize the regime but are not controlled against our base model; our contribution is the sampling strategy and its sample-efficiency benefit, measured by the matched-recipe comparison in Tables~\ref{tab:main}--\ref{tab:curve}.

\subsection{Difficulty Hypothesis Test}
Before changing training, we test the premise that the offline difficulty score predicts where the policy fails. We correlate $\dk$ with the uniform baseline's per-task SR across all $40$ tasks. At a non-saturated checkpoint the correlation is statistically significant (Spearman $\rho=-0.32$, $p<0.05$): higher transition density implies lower success. The relation is clearest at the suite level (the \textsc{Long} suite is highest in $\dk$ and lowest in SR) and is moderate at the task level, limited by ceiling effects on the easy suites; we therefore use a non-saturated checkpoint where the signal is least compressed. A component analysis explains the two-component form of Eq.~\eqref{eq:dk}: the plateau ratio is the dominant predictor on LIBERO, length adds complementary signal, and the gripper-event count is uninformative on this benchmark.

\begin{figure}[t]
\centering
\fbox{\parbox{0.9\columnwidth}{\centering \vspace{2.1cm} TODO: Exp-0 scatter / binned plot of $\dk$ vs.\ baseline per-task SR --- insert \texttt{density\_vs\_sr.png}. \vspace{2.1cm}}}
\caption{Offline transition-density difficulty $\dk$ vs.\ measured per-task success rate (Spearman $\rho=-0.32$, $p<0.05$). Difficulty computed from demonstrations alone predicts where the policy fails.}
\label{fig:exp0}
\end{figure}

\subsection{Ablation Study}
We ablate the components of \tds{} at the $40$k budget; all variants change only the sampling distribution. Table~\ref{tab:ablation} reports (i) a \emph{measured-difficulty} baseline that derives sampling weights from the model's own per-task SR (the ``ask-the-model'' alternative), isolating the value of the \emph{free} data-driven difficulty; (ii) a \emph{uniform-interleaved} control (the \tds{} sampler with uniform weights), isolating the effect of the sampling \emph{mechanism} from the difficulty \emph{weights}; (iii) an \emph{inverted-weights} variant that over-samples easy tasks, testing whether the \emph{direction} of the signal matters or any non-uniform perturbation would help; and (iv) a \emph{no-floor} variant ($\rho{=}0$), testing the uniform floor's role in preventing easy-task forgetting.

\begin{table}[t]
\centering
\setlength{\tabcolsep}{1.2mm}
\begin{tabular}{lccccc}
\toprule
Variant @ $40$k & \textsc{Sp.} & \textsc{Obj.} & \textsc{Goal} & \textsc{Long} & Avg \\
\midrule
Uniform & 73.0 & 95.5 & 38.5 & 40.5 & 61.9 \\
\tds{} (ours) & 97.5 & 98.5 & 92.0 & 71.0 & \textbf{89.75} \\
\midrule
Measured-difficulty & \na{} & \na{} & \na{} & \na{} & \na{} \\
Uniform-interleaved$^\dagger$ & 95.0 & 93.5 & 83.5 & 67.5 & 84.9 \\
Inverted weights & \na{} & \na{} & \na{} & \na{} & \na{} \\
No floor ($\rho{=}0$) & \na{} & \na{} & \na{} & \na{} & \na{} \\
\bottomrule
\end{tabular}
\caption{\textbf{Ablations} (SR \%, $40$k budget). The \emph{uniform-interleaved} control (our two-level sampler with flat weights) decomposes the headline gain: it already reaches $84.9$, so the bulk of the improvement over naive uniform ($61.9$) comes from the sampler's fine-grained batch decorrelation rather than from the difficulty weights. The transition-density weighting itself adds a further $+2.0$ on average measured on a single machine (uniform-interleaved $84.9$ vs.\ \tds{} $86.9$, both on the second machine), concentrated on the up-weighted \textsc{Goal} suite. $^\dagger$Evaluated on the second machine (yyk); the same-machine (sapi) evaluation of this control is pending, so the cross-row comparison to uniform/\tds{} above is approximate ($\sim\!3$ pt machine offset). \textbf{[TODO: same-machine uniform-interleaved + measured-difficulty / inverted / no-floor rows.]}}
\label{tab:ablation}
\end{table}

\paragraph{Decomposing the gain.} The uniform-interleaved control is the most informative ablation: it shows that the headline improvement over the naive uniform baseline has two distinct sources. First, a \emph{mechanism} effect---replacing per-suite episode streaming with fine-grained per-episode interleaving decorrelates the training batch (from $\sim\!4$ to $\sim\!64$ distinct episodes per batch of $64$) and accounts for the majority of the gain ($61.9\!\rightarrow\!84.9$), with the largest lift on the hard suites that the coarse sampler under-mixed. Second, the \emph{transition-density weighting} adds a smaller, targeted increment on top of the well-mixed sampler ($+2.0$ on average same-machine, concentrated on the up-weighted \textsc{Goal} suite). We report this decomposition transparently: the practical recommendation is the full pipeline (interleaving $+$ density weighting), but the novel, model-free component---transition-density reweighting---should be credited with the incremental gain over a properly-mixed uniform baseline, not the full margin over the naive one. \textbf{[TODO: finalize once the same-machine uniform-interleaved and multi-seed runs land; the $+2.0$ weighting effect is currently near the cross-machine / $20$-trial noise floor.]}

\subsection{Generalization and Robustness}
\tds{}'s gains are largest on the suites that stress generalization and horizon: \textsc{Goal} (goal/skill generalization) and \textsc{Long} (long-horizon, multi-stage), where the uniform baseline is weakest (Table~\ref{tab:main}). At $40$k, \tds{} lifts \textsc{Long} from $40.5$ to $71.0$ and \textsc{Goal} from $38.5$ to $92.0$, indicating that prioritizing high-transition-density tasks disproportionately benefits the hardest, longest-horizon behaviors. \textbf{[TODO: cross-embodiment / second-benchmark (e.g., CALVIN) generalization is left to future work; LIBERO is the evaluation here.]}

\subsection{Qualitative Results}
\textbf{[TODO: per-task qualitative panels --- uniform vs.\ \tds{} success rate grouped by difficulty (high/medium/low). On the high-difficulty group, \tds{} is far above uniform; on the low-difficulty group the two are comparable. Insert the generated per-task figure grid.]}

\section{Discussion}
\paragraph{Why TDS works.}
Under uniform sampling, optimization is split evenly across tasks regardless of how much each still has to learn; easy tasks that saturate early continue to consume budget. \tds{} reallocates that budget toward high-transition-density tasks, which are precisely the under-learned, low-success tasks (Fig.~\ref{fig:exp0}). The per-suite gains track the uniform baseline's remaining headroom---near-zero on the already-saturated \textsc{Object} suite, largest on \textsc{Goal}/\textsc{Long}---consistent with this reallocation account.
\paragraph{When it fails.}
The transition signal is computed from demonstrated \emph{actions}; difficulty that stems from \emph{perception} (distractors, lighting, appearance shift) rather than control is invisible to it. \tds{} is therefore complementary to perception-driven or loss-based signals, not a replacement. The constants $(\theta_p,L_p)$ were set on LIBERO and their transfer across control frequencies remains to be verified.
\paragraph{Overhead.}
\tds{} computes its weights once, offline, on CPU in minutes, and the two-level sampler adds negligible data-loading cost and no GPU cost; training throughput is unchanged. Model-coupled alternatives instead require training a reference/proxy model and/or periodic policy evaluation---a multiplicative overhead that \tds{} avoids entirely.

\section{Conclusion}
We presented Transition-Aware Dynamic Sampling, a training strategy that reads per-task manipulation difficulty directly from the transition structure of demonstrations and dynamically reweights task sampling to concentrate learning on hard tasks. As the central contribution of SimVLA, \tds{} is model-free, adds no training overhead, and improves sample efficiency substantially on LIBERO---$+27.9$ average success rate at a fixed low budget and $\sim\!2\times$ fewer steps to match the uniform baseline---while a difficulty score computed before training significantly predicts where the policy fails. The broader message is that demonstration data carries more supervision about \emph{what to train on} than the uniform objective exploits, and that difficulty is better read from the data than asked of the model.

\bibliography{references}

\end{document}
